\chapter{Processamento de Radar Passivo}

A cadeia de radar passivo equaliza os canais de referência e vigilância, cancela o sinal direto e o clutter, calcula a função de ambiguidade cruzada, detecta com controle de falso alarme e associa as detecções entre os nós. A saída é uma trilha com incerteza, não apenas uma coordenada. A detecção de emissor direto deve permanecer separada da detecção por reflexão biestática.

\section{Experimento de consistência de transmissão aeronáutica}
O sinal aeronáutico de referência é o ADS-B, principalmente o 1090ES em 1090 MHz (\texttt{EXP-001}) na aviação comercial. O canal alternativo UAT em 978 MHz (\texttt{EXP-002}) é contextual e relevante sobretudo para determinados usos da aviação geral nos Estados Unidos. Ambos são atribuídos à abertura ou array aeronáutico dedicado de 960--1215 MHz, e não à Yagi UHF de 470--860 MHz. Um receptor diplexado fornece caminhos simultâneos e independentes por serviço, cada qual com preselector, limitador/LNA, ADC coerente e canal FPGA próprios. As mensagens ADS-B são transmitidas automaticamente e repetidamente, com periodicidade dependente do tipo de informação; não constituem um único pacote contínuo contendo todos os parâmetros em todos os instantes. Os campos esperados incluem posição derivada de GNSS, altitude, velocidade, trajetória e identificação da aeronave.

Dois nós radome (\texttt{EXP-005}), separados por uma linha de base proposta de $100\,\mathrm{km}$ (\texttt{EXP-003}), recebem de forma independente a transmissão da aeronave. Essa geometria ilustrativa é distinta do demonstrador-alvo de três nós (\texttt{ARC-002}). Para cada nó, o sistema registra potência recebida $P_i$, ângulo de chegada local $\mathrm{AOA}_i$, timestamp de hardware $t_i$ e estimativa Doppler $f_{D,i}$. A fusão calcula:
\[
\Delta t=t_A-t_B,\qquad \Delta f_D=f_{D,A}-f_{D,B}.
\]
Com as coordenadas levantadas dos nós, hipóteses de efemérides da aeronave e incerteza de propagação, esses observáveis produzem uma região de consistência para altitude, posição e velocidade. O resultado verifica se os parâmetros transmitidos são compatíveis com a observação física; não é uma garantia cega de que todo parâmetro transmitido seja verdadeiro.

O ADS-B fornece uma identidade e um estado cooperativos, enquanto a geometria RF dos dois nós fornece uma verificação independente baseada em potência recebida, AOA, TDOA e Doppler.

A linha de base nominal de 100 km permanece uma geometria de triagem, não uma distância de implantação aprovada. Com as alturas propostas para triagem de 1000 m para a estação e 10000 m para a aeronave (\texttt{EXP-010}), um cálculo padrão de horizonte radioelétrico com Terra efetiva fornece aproximadamente 542 km de horizonte geométrico mútuo. Isso demonstra apenas que a curvatura não exclui 100 km nesse exemplo. Mascaramento pelo terreno, diagramas das antenas, margem de enlace, visibilidade comum, GDOP, autorização de espectro e coordenadas reais dos sítios ainda precisam ser fechados antes do ensaio de campo.

\section{Iluminadores UHF independentes}
Transmissores conhecidos de televisão UHF ou de outros serviços de radiodifusão podem fornecer um teste independente de triangulação junto com o ADS-B. Seus sinais diretos podem ser usados para calibrar AOA, sincronismo, mapeamento de potência recebida e geometria do sítio quando as coordenadas do transmissor e as características da forma de onda forem conhecidas. Seus sinais refletidos também podem apoiar um experimento biestático separado, mas a verificação de emissor direto e a localização por reflexão do alvo não devem ser misturadas no mesmo estimador sem um modelo explícito de propagação.

O experimento deve comparar: (i) ADS-B em 1090 MHz como emissor direto cooperativo; (ii) sinais UHF conhecidos de radiodifusão, provisoriamente na região de 470--860 MHz (\texttt{EXP-004}), como iluminadores diretos de referência; e (iii) ecos biestáticos de um alvo controlado ou rastreado independentemente. A concordância entre esses canais testa a triangulação, a calibração e o relatório de covariância da rede de radomes, em vez de apenas repetir a mensagem da aeronave.

\section{Protocolos experimentais independentes}
O protocolo de emissor direto cooperativo (\texttt{EXP-006}) registra identidade e estado decodificados, potência recebida, AOA, tempo de hardware e Doppler. Sua verdade-terreno é o estado derivado de GNSS e decodificado, com coordenadas dos receptores levantadas independentemente; seu estimador verifica a consistência RF versus mensagem, e suas métricas são taxa de associação e resíduos de posição, velocidade, AOA, tempo e Doppler com incerteza declarada.

O protocolo de calibração por transmissor conhecido (\texttt{EXP-007}) registra a forma de onda UHF direta, potência recebida, AOA e atraso calibrado do canal. Sua verdade-terreno é a posição do transmissor, o setor da antena e os metadados da forma de onda levantados independentemente; seu estimador resolve o caminho direto esperado, e suas métricas são resíduos de AOA, atraso e mapa de potência, além da repetibilidade versus tempo e temperatura.

O protocolo de reflexão biestática (\texttt{EXP-008}) usa canais simultâneos de referência e vigilância, cancelamento do caminho direto e estimador por função de ambiguidade cruzada. Sua verdade-terreno é a trajetória de um alvo controlado ou rastreado independentemente, junto com coordenadas levantadas do transmissor e dos receptores; suas métricas são probabilidade de detecção, probabilidade de falso alarme e erros de atraso, Doppler e localização com consistência de covariância. Nenhum resultado dos dois protocolos de sinal direto substitui essa evidência biestática.

\section{Elipsoide biestático e redução da região de incerteza}
Se a posição do iluminador é $\mathbf{s}$, a posição levantada do receptor $i$ é $\mathbf{r}_i$ e a posição hipotética do alvo é $\mathbf{x}$, o atraso excedente entre o eco e o caminho direto de referência satisfaz
\[
c\,\Delta\tau_i=
\lVert\mathbf{s}-\mathbf{x}\rVert+
\lVert\mathbf{x}-\mathbf{r}_i\rVert-
\lVert\mathbf{s}-\mathbf{r}_i\rVert.
\]
Para $\Delta\tau_i$ fixo, os pontos compatíveis formam uma superfície elipsoidal de isoalcance biestático, com focos no transmissor e no receptor. A potência recebida por face, corrigida pelo diagrama embarcado, perdas e potência efetivamente irradiada do iluminador, fornece uma verossimilhança angular e de enlace; AOA restringe a solução a uma direção ou cone. Observáveis de polarização podem acrescentar discriminação somente quando pertencem à mesma frequência e a portas coerentes calibradas. As Yagis VHF e UHF atuais, por operarem em faixas distintas, não podem ser cruzadas para produzir uma medida polarimétrica do mesmo eco.

A interseção probabilística entre elipsoide de atraso, AOA, Doppler/FDOA, potência e polarização válida produz uma região de estado com covariância, e não um ponto exato. A recepção associada do mesmo eco por múltiplas antenas e faces de um radome acrescenta observações angulares, de potência, atraso e Doppler à estimação. A recepção desse eco por um segundo radome sincronizado acrescenta outro elipsoide e outras restrições; nós integrados adicionais acrescentam informação independente e ajudam a rejeitar ambiguidades.

Em uma reconciliação por mínimos quadrados ponderados ou em um estimador probabilístico equivalente, cada observação válida adicional acrescenta uma matriz de informação positiva semidefinida. Portanto, a covariância estimada não aumenta e a incerteza de medição diminui estritamente nas direções em que a nova observação fornece informação não redundante. Essa propagação conjunta de grandezas de entrada, covariâncias e grandezas de saída multivariadas segue o tratamento metrológico do Suplemento 2 do GUM, JCGM 102:2011 \cite{jcgm1022011}. É uma consequência matemática do modelo de fusão, não uma alegação de desempenho já medido. O conhecimento prévio da posição e da forma de onda do iluminador, os canais de referência, o timestamping em hardware e a base temporal comum, centralmente referenciada e continuamente calibrada por GNSS, referência de frequência, White Rabbit, holdover e medição de atraso permitem associar e reconciliar as observações. O demonstrador deve quantificar a magnitude da redução e confrontar a covariância prevista com erro observado e GDOP; geometria, SNR, multipercurso, calibração e correlação de erros determinam o ganho obtido, não a existência do princípio de redução.

\section{Jamming como emissão observável}
Ataques eletrônicos contra radares ativos podem empregar ruído, varredura, agilidade de frequência e formas de onda coerentes ou enganosas. Para o RADOME, essa energia não desaparece: ela constitui uma emissão direta adicional no espectro e aumenta a detectabilidade passiva do transmissor quando está dentro das faixas e da faixa dinâmica cobertas, sobretudo com potência, largura de banda ou ciclo de trabalho elevados. A literatura da NATO trata separadamente radar passivo, localização passiva de emissores e ataques eletrônicos coerentes contra sensores, inclusive PCL \cite{nato2019passiveradar,nato2017coherentea}.

O protocolo \texttt{EXP-011} trata o jamming como fonte não cooperativa. Múltiplas faces e radomes registram potência, espectro/ocupação, AOA, timestamps e FDOA/Doppler; associação temporal e espectral permite estimar a posição do centro de emissão e, por sequência de estados, seu vetor de movimento. Assim, aumentar a energia radiada para ofuscar radares ativos fornece ao sistema passivo mais energia direta a interceptar e geolocalizar. O processamento deve preservar os dados brutos antes de qualquer supressão de interferência, separar detecção/classificação do jammer da cadeia de ecos biestáticos e usar limitadores, preseleção e controle de ganho para evitar saturação.

Essa vantagem não equivale a identificar automaticamente a aeronave. Jammers direcionais, de baixa probabilidade de interceptação, coerentes, rebocados, descartáveis ou stand-off podem produzir múltiplas fontes, falsos ângulos e associação ambígua; o emissor localizado pode não estar no caça. A afirmação controlada é que o jamming aumenta a observabilidade eletromagnética de sua fonte para a rede passiva e cria observáveis adicionais. A campanha deve medir probabilidade de interceptação, erro de AOA/TDOA/FDOA, saturação, classificação e consistência da trajetória para cada classe de sinal.


\section{Torres celulares como iluminadores calibrados}
As torres celulares acrescentam um conjunto particularmente valioso de iluminadores conhecidos. Suas coordenadas levantadas, orientações setoriais e alocações nominais de portadoras permitem usar múltiplas frequências e formas de onda celulares como referências diretas. O sinal direto pode calibrar cobertura, AOA, sincronismo e consistência de potência recebida; as reflexões em aeronaves, terreno ou alvos controlados podem então testar a detecção biestática em várias portadoras e larguras de banda. Cada portadora deve ser tratada como um caso separado de forma de onda e propagação, registrando potência do transmissor, setor da antena, agendamento, sincronismo e RFI local nos metadados.

Isso transforma a infraestrutura celular em um campo de testes distribuído para o conceito de radar de múltiplos radomes de banda larga. O sistema não se torna um radar ativo: os radomes continuam sendo receptores passivos, e os observáveis úteis vêm das transmissões conhecidas, da geometria e de seus componentes refletidos.

\begin{figure}[H]
\centering
\includegraphics[width=0.98\textwidth]{pt-BR/fig15_aircraft_two_radomes_100km.png}
\caption{Cenário aeronáutico de dois nós: os dois radomes ampliados definem a linha de base nominal de 100 km, a aeronave está acima e deslocada lateralmente, e os caminhos e cones azul/vermelho representam sinais recebidos e observações angulares independentes usados na consistência de potência, AOA, TDOA e Doppler.}
\label{fig:fig15_aircraft_two_radomes_100km}
\end{figure}

\begin{figure}[H]
\centering
\includegraphics[width=0.95\textwidth]{pt-BR/fig09_fluxo_processamento.png}
\caption{Fluxo de processamento passivo multiestático: as entradas de referência e vigilância convergem para equalização, cancelamento do caminho direto, processamento atraso--Doppler, detecção controlada e associação; os ramos inferiores representam prioris de emissores e a estimativa de estado resultante.}
\label{fig:fig09_fluxo_processamento}
\end{figure}
